\documentclass[11pt]{beamer}
\usepackage[utf8]{inputenc}
\usepackage[T1]{fontenc}
\usepackage[slovak]{babel}
\usepackage{listings}
\usepackage{xcolor}

\usetheme{Madrid}
\usecolortheme{default}

\title{Kompilátor Datalogu s funkčnými symbolmi (a vstavanými operátormi) do relačnej algebry}
\author{Matúš Duchyňa}
\institute{
    Fakulta matematiky, fyziky a informatiky \\
    Univerzita Komenského v Bratislave
    \vspace{0.5em} % Pridanie malého priestoru
    \\
    Vedúci práce: doc. Mgr. Tomáš Plachetka, Dr.
}
\date{\today}
    
\newcommand{\customfootertitle}{Kompilátor Datalog s funkčnými symbolmi -> Relačná algebra}

\setbeamertemplate{footline}{%
    \leavevmode%
    \begin{beamercolorbox}[wd=\paperwidth,sep=0.4em]{section in foot}%
        \insertshortauthor \hspace*{0.5em} | \hspace*{0.5em}
        \customfootertitle
        \hfill%
        \insertshortdate \hspace*{0.5em} | \hspace*{0.5em}
        \insertframenumber/\inserttotalframenumber\hspace*{1em}
    \end{beamercolorbox}%
}

\lstset{
    basicstyle=\ttfamily\small,
    breaklines=true,
    frame=single,
    language=Java,
    showstringspaces=false,
    commentstyle=\color{gray},
    keywordstyle=\color{blue},
}

\begin{document}

\frame{\titlepage}

\begin{frame}
\frametitle{Obsah práce}
\framesubtitle{Dalalog}
\begin{itemize}
    \item Definovať datalog s funkčnými symbolmi
    \item Definovať syntax
    \item Spraviť gramatiku pre ANTLR
    \item Spraviť kompilátor do relačnej algebry
\end{itemize}
\end{frame}

\begin{frame}
\frametitle{Obsah práce}
\framesubtitle{Relačná algebra (nadstavba na minulý rok)}
\begin{itemize}
    \item Identifikovať aké operácie potrebujeme
    \item Definovať syntax
\end{itemize}
\end{frame}

\begin{frame}
\frametitle{Obsah práce}
\framesubtitle{Kód}
Pri dobre zvládnutých predošlých krokoch priamočiara robota
\begin{itemize}
    \item Samotný kompilátor
    \item Tester lexera
    \item Tester parsera
    \item Overiť funkčnosť
\end{itemize}
\end{frame}

\begin{frame}
\frametitle{Obsah práce}
\framesubtitle{Kapitoly}
% TODO datumy

\begin{enumerate}
    \item Datalog \qquad 25.3.
    \begin{enumerate} 
        \item Definície \quad Ešte zopár viet mi tam chýba
        \item Rozšírenie datalogu o funkčné symboly \quad Done
        \item Syntax rozšíreného datalogu \quad Done
        \item Well-founded model \quad Napísať
    \end{enumerate}
    \item Relačná algebra \qquad 1.4. 
    \begin{enumerate}
        \item Definície \quad Done
        \item Inklúzia medzi učebnicovou a našou relačnou algebrou \quad Done
    \end{enumerate}
    \item Kompilátor \qquad 15.4. (program funkčný)
    \begin{enumerate}
        \item ANTRL \quad Napísať
        \item Gramatika datalogu pre ANTLR \quad Napísať
        \item Algoritmus prekladu \quad Napísať
    \end{enumerate}
    \item Kontroly \qquad 29. 4.
\end{enumerate}
\end{frame}

\end{document}